\subsection{Trace of an endomorphism}

Let E and F be two hilbertian spaces. Conforming to the conventions of V, p.~40, we let \(ba^{*}\) , for \(a\) in E and \(b\) in F, denote the continuous linear mapping \(x \mapsto b \langle a|x\rangle\) from E into F.

Lemma 1. --- There exists an isomorphism \(\theta\) from the vector space \(\mathbf{F} \otimes \mathbf{E}'\) onto the space \(\mathcal{L}_f(\mathbf{E}; \mathbf{F})\) of all finite rank continuous linear mappings from \(\mathbf{E}\) into \(\mathbf{F}\) , characterized by \(\theta(b \otimes a^*) = ba^*\) for \(a \in \mathbf{E}\) , \(b \in \mathbf{F}\) .

By A, II, § 4, No.~2, there exists an injective linear mapping \(\theta\) from \(\mathbf{F} \otimes \mathbf{E}'\) into \(\mathcal{L}(\mathbf{E};\mathbf{F})\) and only one such, which transforms \(b \otimes a'\) into the linear mapping \(x \mapsto ba'(x)\) for \(a' \in \mathbf{E}'\) , \(b \in \mathbf{F}\) . Evidently \(\theta(b \otimes a^*) = ba^*\) , and the image of \(\theta\) is contained in \(\mathcal{L}_f(\mathbf{E};\mathbf{F})\) . However, let \(u \in \mathcal{L}_f(\mathbf{E};\mathbf{F})\) and let \((e_1,\dots,e_n)\) be an orthonormal basis of the image of \(u\) in \(\mathbf{F}\) . Let \(f_i = u^*(e_i)\) for \(1 \leqslant i \leqslant n\) . For every \(x \in \mathbf{E}\) , we have

\[
u (x) = \sum_ {i = 1} ^ {n} \left\langle e _ {i} | u (x) \right\rangle . e _ {i} = \sum_ {i = 1} ^ {n} \left\langle f _ {i} | x \right\rangle . e _ {i},
\]

hence \(u = \sum_{i=1}^{n} e_i f_i^* = \theta(\sum_{i=1}^{n} e_i \otimes f_i^*)\) . Therefore the image of \(\theta\) is equal to \(\mathcal{L}_f(\mathrm{E};\mathrm{F})\) .

\[
\mathrm{E} = \mathrm{F},
\]

\[
\mathscr {L} _ {f} (\mathrm{E}) = \mathscr {L} _ {f} (\mathrm{E}; \mathrm{E})
\]

\[
\mathcal {L} _ {f} (\mathrm{E})
\]

\[
\tau (\theta (a \otimes a ^ {\prime})) = a ^ {\prime} (a)
\]

\[
a \in \mathrm{E}, a ^ {\prime} \in \mathrm{E} ^ {\prime}\tag{22}
\]

\[
\tau (b a ^ {*}) = \langle a | b \rangle \quad \text { for } a, b \text { in } E.
\]

When E is finite dimensional, we have \(\mathcal{L}_{f}(\mathrm{E}) = \mathcal{L}(\mathrm{E})\) and \(\tau(u)\) is the trace of the endomorphism u of E (A, II, § 4, No.~3).

Lemma 2. --- Let \((e_i)_{i\in \mathbf{I}}\) be an orthonormal basis of E. Then

\[
\tau (u) = \sum_ {i \in \mathrm{I}} \left\langle e _ {i} | u (e _ {i}) \right\rangle
\]

for all \(u \in \mathcal{L}_f(\mathrm{E})\) .

It is enough to consider the case when \(u = ba^{*}\) with \(a, b\) in E. Then

\[
\langle e _ {i} | u (e _ {i}) \rangle = e _ {i} ^ {*} b. a ^ {*} e _ {i} = \overline {{\langle e _ {i} | a \rangle}} \langle e _ {i} | b \rangle
\]

and lemma 2 follows from formula (22) and formula (3) of V, p.~22.

Lemma 3. --- Let u be a continuous and positive endomorphism of E, and F the set of all finite rank orthoprojectors on E. Then for every orthonormal basis \((e_{i})_{i\in\mathbf{I}}\) of E, we have \((in\overline{\mathbf{R}}_{+})\) the equality

\[
\sum_ {i \in \mathrm{I}} \left\langle e _ {i} | u (e _ {i}) \right\rangle = \sup _ {p \in \mathcal {F}} \tau (p u p).
\]

For every finite subset J of I, put \(p_{J} = \sum_{i \in J} e_{i} e_{i}^{*}\) ; this is the orthoprojector from E onto the vector subspace generated by the vectors \(e_{i}\) , where i ranges over J. We have

\[
p _ {\mathrm{J}} u p _ {\mathrm{J}} = \sum_ {i \in \mathrm{J}, j \in \mathrm{J}} \left\langle e _ {i} | u (e _ {j}) \right\rangle e _ {i} e _ {j} ^ {*},
\]

hence \(\tau(p_{\mathbf{J}}up_{\mathbf{J}}) = \sum_{i\in \mathbf{J}}\langle e_i|u(e_i)\rangle\) . Since \(p_{\mathbf{J}}\in \mathcal{F}\) ,

\[
\sum_ {i \in \mathbf {J}} \left\langle e _ {i} | u (e _ {i}) \right\rangle \leqslant \sup _ {p \in \mathcal {F}} \tau (p u p);
\]

and so we conclude that

\[
\sum_ {i \in \mathbf {I}} \left\langle e _ {i} | u (e _ {i}) \right\rangle = \sup _ {\mathbf {J}} \sum_ {i \in \mathbf {J}} \left\langle e _ {i} | u (e _ {i}) \right\rangle \leqslant \sup _ {p \in \mathcal {F}} \tau (p u p).
\]

Let \(v\) be a finite rank continuous and positive endomorphism of \(\mathbf{E}\) and let \(p \in \mathcal{F}\) . By th. 2 of \(\mathbf{V}\) , p.~23 there exists an orthonormal basis \((f_{\alpha})_{\alpha \in \mathbf{A}}\) of \(\mathbf{E}\) and a finite subset \(\mathbf{B}\) of \(\mathbf{A}\) such that \((f_{\alpha})_{\alpha \in \mathbf{B}}\) is an orthonormal basis of the image of \(p\) . Then we have \(p = \sum_{\alpha \in \mathbf{B}} f_{\alpha} f_{\alpha}^{*}\) , and so, as above, the relation \(\tau(pvp) = \sum_{\alpha \in \mathbf{B}} \langle f_{\alpha}|v(f_{\alpha}) \rangle\) . By lemma 2 (V, p.~48) we have \(\tau(v) = \sum_{\alpha \in \mathbf{A}} \langle f_{\alpha}|v(f_{\alpha}) \rangle\) , which gives the formula

\[
\sum_ {\alpha \in \mathbf {B}} \langle f _ {\alpha} | v (f _ {\alpha}) \rangle \leqslant \tau (v).
\]

Applying this inequality to the case where \(v = p_{\mathbf{J}}.up_{\mathbf{J}}\) and where \(\mathbf{J}\) is a finite subset of \(\mathbf{I}\) , we get

\[
\sum_ {\alpha \in \mathbf {B}} \left\langle p _ {\mathbf {J}} (f _ {\alpha}) | u p _ {\mathbf {J}} (f _ {\alpha}) \right\rangle \leqslant \sum_ {i \in \mathbf {J}} \left\langle e _ {i} | u (e _ {i}) \right\rangle .\tag{23}
\]

For every \(x \in \mathbb{E}\) , we have \(p_{\mathbf{J}}(x) = \sum_{i \in \mathbf{J}} \langle e_i | x \rangle e_i\) , and so \(x = \lim_{\mathbf{J}} p_{\mathbf{J}}(x)\) with respect to the ordered directed set of finite subsets \(\mathbf{J}\) of \(\mathbf{I}\) . Passing to the limit over \(\mathbf{J}\) in (23), we get

\[
\tau (p u p) = \sum_ {\alpha \in \mathrm{B}} \left\langle f _ {\alpha} | u (f _ {\alpha}) \right\rangle \leqslant \sum_ {i \in \mathrm{I}} \left\langle e _ {i} | u (e _ {i}) \right\rangle ,
\]

and this completes the proof of lemma 3.

DEFINITION 7. --- Let u be a continuous and positive endomorphism of the hilbertian space E. Let

\[
\operatorname{Tr} (u) = \sup _ {p \in \mathcal {F}} \tau (p u p)\tag{24}
\]

(upper bound in \(\overline{\mathbf{R}}_+\) ), where \(\mathcal{F}\) is the set of all finite rank orthoprojectors on E. We say that \(\operatorname{Tr}(u)\) is the trace of \(u\) .

Let \(p\) be the orthoprojector from \(E\) onto a finite dimensional vector subspace of \(E\) , and let \((x_1, \ldots, x_m)\) be an orthonormal basis of \(F\) . We have established the relation \(\tau(pup) = \sum_{i=1}^{m} \langle x_i | u(x_i) \rangle\) . Consequently, we can define the trace by the formula

\[
\operatorname{Tr} (u) = \sup _ {x _ {1}, \dots , x _ {m}} \sum_ {i = 1} ^ {m} \left\langle x _ {i} | u (x _ {i}) \right\rangle ,\tag{24'}
\]

where \((x_{1},\dots,x_{m})\) ranges over the set of all finite orthonormal sequences of vectors of E.

By lemma 3 (V, p.~48), we have

\[
\operatorname{Tr} (u) = \sum_ {i \in \mathbf {I}} \left\langle e _ {i} | u (e _ {i}) \right\rangle\tag{25}
\]

for every orthonormal basis \((e_{i})_{i\in I}\) of E. From this, we deduce

\[
\operatorname{Tr} (u + v) = \operatorname{Tr} (u) + \operatorname{Tr} (v)\tag{26}
\]

\[
\operatorname{Tr} (\lambda u) = \lambda . \operatorname{Tr} (u)\tag{27}
\]

for all continuous and positive endomorphisms u and v of E and for every real number \(\lambda \geqslant 0\) (we make the convention \(0. (+\infty) = 0\) in (27)). Let \(\phi\) be an isomorphism from E onto a hilbertian space F; since \(\phi\) transforms every orthonormal basis of E into an orthonormal basis of F, we get from (25) that

\[
\operatorname{Tr} \left(\phi u \phi^ {- 1}\right) = \operatorname{Tr} (u).\tag{28}
\]

Let \((u_{\alpha})_{\alpha \in \mathbf{A}}\) be a non-empty directed increasing and bounded family of continuous and positive endomorphisms of E; let \(u = \sup_{\alpha} u_{\alpha}\) , then \(\langle x|u(x) \rangle = \sup_{\alpha} \langle x|u_{\alpha}(x) \rangle\) for all \(x \in \mathrm{E}\) (V, p.~46, prop. 13). We have \(\operatorname{Tr}(u) = \sup_{\mathbf{J} \subset \mathbf{I}} \sum_{i \in \mathbf{J}} \langle e_i | u(e_i) \rangle\) , where \(\mathbf{J}\) ranges over all finite subsets of I, hence

\[
\operatorname{Tr} (u) = \sup _ {\alpha} \operatorname{Tr} (u _ {\alpha}) \quad \text { for } \quad u = \sup _ {\alpha} u _ {\alpha}.\tag{29}
\]

Let \(p_{\mathrm{F}}\) be the orthoprojector from E onto the hilbertian subspace F; there exists an orthonormal basis \((e_i)_{i \in \mathrm{I}}\) of E and a subset J of I, such that \((e_i)_{i \in \mathrm{J}}\) is an orthonormal basis of F. We have \(\operatorname{Tr}(p_{\mathrm{F}} up_{\mathrm{F}}) = \sum_{i \in \mathrm{J}} \langle e_i | u(e_i) \rangle\) . This formula has two consequences: firstly, we have \(\operatorname{Tr}(p_{\mathrm{F}} up_{\mathrm{F}}) \leqslant \operatorname{Tr}(u)\) ; secondly, taking \(u = 1_{\mathrm{E}}\) , we get

\[
\operatorname{Tr} (p _ {\mathrm{F}}) = \left\{ \begin{array}{l l} \dim \mathrm{F} & \text { if   F   is   finite   dimensional } \\ + \infty & \text { if   not. } \end{array} \right.\tag{30}
\]

DEFINITION 8. --- Let E be a complex hilbertian space. We write \(\mathcal{L}^{1}(\mathrm{E})\) for the vector subspace of \(\mathcal{L}(\mathrm{E})\) generated by all continuous, positive endomorphisms of E with finite trace.

By formula (25) of V, p.~50, the trace extends to a linear form on \(\mathcal{L}^1 (\mathrm{E})\) , again denoted by Tr, and satisfying the relation \(\operatorname{Tr}(u) = \sum_{i\in I}\langle e_i|u(e_i)\rangle\) for all \(u\) in \(\mathcal{L}^1 (\mathrm{E})\) and for every orthonormal basis \((e_i)_{i\in I}\) of E. For every \(u\in \mathcal{L}^1 (\mathrm{E})\) , we have \(u^{*}\in \mathcal{L}^{1}(\mathrm{E})\) and \(\operatorname {Tr}(u^{*}) = \overline{\operatorname{Tr}(u)}\) . Formula (28) of V, p.~50 extends to the case where \(u\) belongs to \(\mathcal{L}^1 (\mathrm{E})\) . Let F be a hilbertian subspace of E; by formula (30), the orthoprojector \(p_{\mathrm{F}}\) belongs to \(\mathcal{L}^1 (\mathrm{E})\) if and only if F is finite dimensional. For every \(a\) and \(b\) in E, we have \(4ab^{*} = \sum_{\varepsilon^{4} = 1}\varepsilon (a + \varepsilon b)(a + \varepsilon b)^{*}\) and \(cc^{*}\) is a positive operator with finite trace for all \(c\in \mathrm{E}\) : consequently, if \(u\) is a finite rank, continuous endomorphism of E, then \(u \in \mathcal{L}^1(\mathrm{E})\) and \(\operatorname{Tr}(u) = \tau(u)\) .

Let \(\mathbf{E}\) be a real hilbertian space, and let \(\mathrm{E}_{(\mathbf{C})}\) be its complexification (V, p.~5). We identify \(\mathbf{E}\) with a subset of \(\mathrm{E}_{(\mathbf{C})}\) . Then \(\mathcal{L}(\mathrm{E})\) can be identified with a real vector subspace of \(\mathcal{L}(\mathrm{E}_{(\mathbf{C})})\) consisting of all continuous linear mappings \(u\) from \(\mathrm{E}_{(\mathbf{C})}\) into \(\mathrm{E}_{(\mathbf{C})}\) such that \(u(\mathrm{E}) \subset \mathrm{E}\) . In this case we write \(\mathcal{L}^1(\mathrm{E}) = \mathcal{L}(\mathrm{E}) \cap \mathcal{L}^1(\mathrm{E}_{(\mathbf{C})})\) . For every \(u \in \mathcal{L}^1(\mathrm{E})\) , the trace \(\operatorname{Tr}(u)\) is real and is equal to \(\operatorname{Tr}(u^*)\) . Formulas (25) and (28) are again valid, \(\mathcal{L}_f(\mathrm{E}) \subset \mathcal{L}^1(\mathrm{E})\) and \(\operatorname{Tr}(u) = \tau(u)\) for all \(u \in \mathcal{L}_f(\mathrm{E})\) . Finally, a closed vector subspace \(\mathbf{F}\) of \(\mathbf{E}\) is finite dimensional if and only if \(p_{\mathrm{F}}\) belongs to \(\mathcal{L}^1(\mathrm{E})\) .

* Remark 1. --- We shall later define the notion of a nuclear mapping from a Banach space E into a Banach space F. We shall show that when \(\mathcal{L}^1 (\mathrm{E})\) consists of all nuclear mappings from E into E, then E is a real or complex hilbertian space. *

PROPOSITION 14. --- Let \(E_{1}, \ldots, E_{n}\) be hilbertian spaces, \(E = E_{1} \hat{\otimes}_{2} \ldots \hat{\otimes}_{2} E_{n}\) , and \(u_{i}\) a continuous endomorphism of \(E_{i}\) for \(1 \leqslant i \leqslant n\) . If \(u_{1}, \ldots, u_{n}\) are positive, then so is \(u = u_{1} \hat{\otimes}_{2} \ldots \hat{\otimes}_{2} u_{n}\) , and

\[
\operatorname{Tr} (u) = \prod_ {i = 1} ^ {n} \operatorname{Tr} (u _ {i}).\tag{31}
\]

If \(u_{i} \in \mathcal{L}^{1}(\mathrm{E}_{i})\) for all \(1 \leqslant i \leqslant n\) , then \(u \in \mathcal{L}^{1}(\mathrm{E})\) and formula (31) is again valid in this case.

Proceeding by induction on \(n\) , we immediately reduce to the case \(n = 2\) .

For i=1,2, we define a sesquilinear form \(\Phi_{i}\) on \(E_{i}\) by the formula \(\Phi_{i}(x,y)=\langle x|u_{i}(y)\rangle\) for x, y in \(E_{i}\) . If \(u_{1}\) and \(u_{2}\) are positive, the forms \(\Phi_{1}\) and \(\Phi_{2}\) are hermitian and positive. By prop. 1 of V, p.~25 there exists a positive hermitian form \(\Phi\) on the vector space \(E_{1}\otimes E_{2}\) such that

\[
\Phi (x _ {1} \otimes x _ {2}, y _ {1} \otimes y _ {2}) = \Phi_ {1} (x _ {1}, y _ {1}). \Phi_ {2} (x _ {2}, y _ {2})
\]

for \(x_{1}, y_{1}\) in \(E_{1}\) and \(x_{2}, y_{2}\) in \(E_{2}\) . We verify immediately the relation \(\Phi(z, t) = \langle z|u(t) \rangle\) for \(z\) and \(t\) in \(E_{1} \otimes E_{2}\) . Since \(\Phi\) is positive, we have \(\langle z|u(z) \rangle \geqslant 0\) for all \(z\) in \(E_{1} \otimes E_{2}\) . Since \(u\) is continuous and \(E_{1} \otimes E_{2}\) is dense in the hilbertian space \(E = E_{1} \hat{\otimes}_{2} E_{2}\) , we conclude that \(u\) is a continuous and positive endomorphism of \(E\) .

Let \((e_i)_{i\in I}\) be an orthonormal basis of \(\mathbf{E}_1\) and \((f_j)_{j\in J}\) an orthonormal basis of \(\mathbf{E}_2\) ; then the family \((e_i\otimes f_i)_{i\in I,j\in J}\) is an orthonormal basis of E and we have

\[
\begin{array}{r l} & {\mathrm{Tr} (u) = \sum_ {i \in \mathbf {I}} \sum_ {j \in \mathbf {J}} \langle e _ {i} \otimes f _ {j} | u (e _ {i} \otimes f _ {j}) \rangle} \\ & {\quad = \sum_ {i \in \mathbf {I}} \sum_ {j \in \mathbf {J}} \langle e _ {i} | u _ {1} (e _ {i}) \rangle . \langle f _ {j} | u _ {2} (f _ {j}) \rangle} \\ & {\quad = \mathrm{Tr} (u _ {1}). \mathrm{Tr} (u _ {2}).} \end{array}
\]

In particular, if \(u_{1}\) and \(u_{2}\) are positive endomorphisms with finite trace, then so is u. By linearity, we deduce that u belongs to \(\mathcal{L}^{1}(\mathrm{E})\) when K = C and that the \(u_{i}\) belong to \(\mathcal{L}^{1}(\mathrm{E}_{i})\) for i = 1, 2; formula (31) extends to this case by linearity. Finally, the case when K = R and the \(u_{i} \in \mathcal{L}^{1}(\mathrm{E}_{i})\) reduces to the complex case by extension of the scalars.

Remark 2. --- Let E be a hilbertian space, which is the hilbertian sum of a family \((\mathrm{E}_{i})_{i\in\mathrm{I}}\) of hilbertian subspaces. Let u be an element of \(\mathcal{L}(\mathrm{E})\) such that \(u(\mathrm{E}_{i}) \subset \mathrm{E}_{i}\) for all \(i \in I\) ; let \(u_{i}\) be the element of \(\mathcal{L}(\mathrm{E}_{i})\) which coincides with u on \(E_{i}\) . Then \(\operatorname{Tr}(u) = \sum_{i \in \mathrm{I}} \operatorname{Tr}(u_{i})\) when u is positive, or belongs to \(\mathcal{L}^{1}(\mathrm{E})\) ; this relation follows from formula (25) of V, p.~50 applied to an orthonormal basis of E which is the union of orthonormal bases of each of the \(E_{i}\) .

